\section{现代系统延伸：从 Cicero 读 Agent 架构}

本节是 2026 年的工程解读，不属于 2023 年课堂事实。它不引入后来产品作为讲者论据，而是把 Cicero 已经展示的接口翻译成更一般的 agent 设计语言。

\subsection{把“说什么”与“做什么”分开}

现代 agent 常同时维护 private plan、tool actions 与 user-facing response。Cicero 的经验说明，直接把完整内部计划交给语言模型自由表达会产生泄密、矛盾和不可审计问题；更好的做法是先形成结构化 action intent，再由受约束的 communication layer 决定对外信息。

\begin{knowledgebox}{可迁移的五层接口}
\begin{enumerate}
  \item \textbf{State estimator}：把环境、历史和反馈转为当前 belief；
  \item \textbf{Policy prior}：给出常见、低成本、与人类惯例兼容的候选；
  \item \textbf{Planner / search}：按任务价值分配额外 inference compute；
  \item \textbf{Intent interface}：把内部决策压缩成可验证的结构；
  \item \textbf{Communication evaluator}：预测输出会怎样改变环境与他人行为。
\end{enumerate}
\end{knowledgebox}

\subsection{Prompt injection 与战略消息的共同结构}

Diplomacy 对手的消息不是可信系统指令，而是环境中一个带目标的 observation。这与工具型 agent 面对不可信网页、邮件或第三方文档时的安全边界相似：内容可以更新 belief，却不应自动获得控制权。Cicero 的“预测对方行为，再由自身 value 判断”提供了比直接 instruction-following 更稳健的范式。

\begin{warningbox}{类比的边界}
Diplomacy 有封闭规则、可模拟动作和明确计分；开放世界 agent 往往缺少可靠 transition model 与 reward。因而不能简单复制 piKL 或 value filter，而应继承其审计思想：区分数据与指令、显式记录 intent、让高风险输出经过独立评价。
\end{warningbox}

\subsection{成本治理：不是所有请求都值得 search}

推理时计算应有预算策略。一个实用系统需要估计问题价值、验证难度、失败代价和剩余 latency，再决定候选数、搜索深度、tool calls 与 evaluator 强度。若没有停止条件，agent 会把“多想”退化为无限循环；若所有请求固定一次 greedy decode，又放弃了 search 最有价值的动态资源分配。

\begin{importantbox}{工程上的预算问题}
每增加一次 rollout 或 evaluator 调用，都应能回答：它降低了哪类错误，是否有独立验证信号，边际收益何时转负，以及错误相关性会不会让多数投票失效。Inference compute 只有进入可测量的决策流程，才是系统能力，而不是账单增长。
\end{importantbox}

\subsection{本章小结}

Cicero 对现代 agent 最有价值的遗产不是某个特定模型，而是模块之间的契约：不可信语言先进入 belief，不直接控制行动；planning 输出结构化 intent；对外语言在发送前接受行为后果评估；计算预算随任务价值变化。

